\section{Discussion}
\label{sec:discussion}

The shared-checkpoint comparison answers RQ1 at the level of implemented diagnostics. Increasing the candidate budget and applying symbolic scoring improved pc-set, interval-vector, cyclic row-order, rhythmic-profile, and gesture endpoints without changing teacher-forced token accuracy. The contrast therefore concerns the combined guided-decoding procedure rather than a stronger trained language model or reranking alone. The held-out density curve also changed slightly in the preferred direction, although it was not available to the selector and no inferential uncertainty is reported.

The condition-removal experiments address RQ2 more directly. Each removal produced its largest interpretable loss in the associated control dimension. The pronounced pc-set and interval-vector changes indicate that pitch-class tokens do more than label a fixed generator distribution. Likewise, row, rhythm, and gesture fields contribute information that cannot be recovered fully from the remaining prefix. Because the corpus is procedural and each comparison uses one fixed seed, these effects should be understood as internal validity checks rather than evidence of broad musical generalization.

The serial results expose an important distinction between aggregate presence and ordered realization. The proposed decoder covered almost the full chromatic aggregate on serial targets, yet exact serial-transformation accuracy was zero and cyclic row-order accuracy was only 0.2411. Candidate reranking therefore selected fragments with better local or cyclic agreement without solving complete $P/R/I/RI$ realization. The procedural reference's near-perfect strict-form score confirms that the metric is attainable under direct construction and that the remaining gap lies in the learned generation path.

RQ3 is supported only at the structural notation level. Every attempted MusicXML file parsed and matched requested measures and parts, but mean content-span ratio for the proposed model was 0.5195. Measure-count adherence partly reflects explicit padding rather than continuous musical activity. The exports are consequently usable as editable sketches and audit objects, while publication-quality engraving and compositional usefulness remain questions for expert review.

More broadly, the results favor a division of labor between probabilistic generation and explicit symbolic analysis. The neural model supplies variation, while deterministic theory functions expose and rank constraint adherence. The rule reference remains stronger for exact serial and rhythmic realization, suggesting that future assistants may benefit from hybrid decoding in which selected constraints are constructed exactly and neural generation operates within the remaining degrees of freedom.
